\section{Application: Matrix completion}
\label{sec:matrix_completion}

Assume we have some matrix $Y \in \R^{n \times m}$ 
which has low ``complexity.'' 
Perhaps it is low-rank or simple by some other measure.
If we observe $Y_{ij}$ for a limited set of entries
$(i,j) \in E \subset [n] \times [m]$,
then {\it matrix completion} is any method which
constructs a matrix $\hat{Y}$
so that $\| \hat{Y} - Y\|$ is small, or even zero.
Matrix completion has attracted significant attention in recent years as a tractable
algorithm for making recommendations to users of online systems based on the
tastes of other users (a.k.a.\ the Netflix problem).
We can think of it as the matrix version of 
compressed sensing (\citet*{candes2010, candes2010a}).

Recently, a number of authors have studied the performance of 
matrix completion algorithms where the index set $E$ is the edge
set of a regular random graph 
(\citet*{heiman2014, bhojanapalli2014,gamarnik2017}).
\citet*{heiman2014} describe a 
{\it deterministic} method of matrix completion,
where they can give performance guarantees for a fixed observation set $E$ 
over many input matrices $Y$.
The error of their reconstruction depends on the spectral gap of the graph.
We expand upon the result of \citet*{heiman2014}, 
extending it to rectangular matrix and improving their bounds in the process.

\subsection{Matrix norms as measures of complexity and their relationships}

We will employ a number of different matrix and vector norms in this Section.
These are all related by the properties of the underlying Banach spaces.
The complexity of $Y$ 
is measured using a factorization norm (also called the max-norm):
\[
	\gamma_2 (Y) = \min_{UV^* = Y} \|U\|_{\ell_2 \to \ell_\infty^n} \|V\|_{\ell_2 \to \ell_\infty^m}.
\]
The minimum is taken over all possible factorizations of $Y = UV^*$, 
and the norm 
% $\| \cdot \|_{\ell_2 \to \ell_\infty^n}$
$\| X \|_{\ell_2 \to \ell_\infty^n} = \max_i \sqrt{\sum_j X^2_{ij}}$
returns the largest $\ell_2$ norm of a row.
So, equivalently,
\[
	\gamma_2 (Y) = \min_{U V^* = Y} \max_{i,j} \|u_i \|_2 \, \|v_j \|_2 ,
\]
where $u_i$ and $v_i$ are the rows of $U$ and $V$.
See \citet*{linial2007} for a number of results about the norm $\gamma_2$.
In particular, note that
\begin{align}
\gamma_2(A \circ B) \leq \gamma_2(A) \gamma_2(B) \label{eq:submultiplicative} \\
  \label{eq:trace_norm_equiv}
  \frac{1}{\sqrt{n m}} \| Y \|_{\rm Tr} \leq \gamma_2(Y) \\ % \leq \| Y \|_{\rm Tr} \\
  \label{eq:rank_equiv}
  \gamma_2(Y) \leq \sqrt{\mathrm{rank} (Y) } \| Y \|_{\infty}.
\end{align}
Property \eqref{eq:submultiplicative} says that $\gamma_2$
is sub-multiplicative under the Hadamard product \citep*{lee2008,heiman2014} and will be used in our proof.
Properties \eqref{eq:trace_norm_equiv} and \eqref{eq:rank_equiv}
relate $\gamma_2$ to two common complexity measures of matrices,
the trace norm (sum of singular values, i.e.\ the $\ell^m_2 \to \ell^n_2$ nuclear norm) 
and rank.
Note also the well-known fact that
\[
\| Y \|_{\rm Tr} = \min_{UV^* = Y} \| U \|_F \|V \|_F ,
\]
where 
% $\| X \|_F = \sqrt{\sum_{i=1}^n \sum_{j=1}^m X_{ij}^2}$ 
$\| X \|_F = \sqrt{\sum_{ij} X_{ij}^2}$ 
is the Frobenius norm.
We see that the trace norm constrains factors $U$ and $V$ to be small on average
via $\| \cdot \|_F$,
whereas the norm $\gamma_2$ is similar but constrains factors uniformly
via $\| \cdot \|_{\ell_2 \to \ell_\infty^n}$.
However, we should note that computing $\gamma_2(Y)$ is more costly than the
trace norm, which can be performed with just the singular value decomposition,
although still possible in polynomial time with convex programming \citep*{heiman2014}.

% There is one last norm we must introduce.
% \todo{talk about ``nuclear norm'' $\nu$}

\subsection{Matrix completion generalization bounds}

The method of matrix completion that we study
is to return the matrix $X$ which is the solution to:
\begin{equation}
\label{matrix_completion}
\begin{aligned}
& \underset{X}{\text{minimize}}
& & \gamma_2(X) \\
& \text{subject to}
& & X_{ij} = Y_{ij}, \; (i,j) \in E .
\end{aligned}
\end{equation}
The $\gamma_2$ norm was first proposed by
\citet*{srebro2005,srebro2005a} 
as a robust complexity measure,
and it was shown to be an effective and practical regularization
on real datasets \citep*{lee2010a,recht2013}.

\citet*{heiman2014} analyze the performance
of the convex program \eqref{matrix_completion} 
for a square matrix $Y$
using an expander argument,
assuming that $E$ is the edge set of a $d$-regular graph 
with second eigenvalue $\eta$.
They obtain the following theorem:
\begin{theorem}[\citet*{heiman2014}]
\label{thm:square_MSE}
Let $E$ be the set of edges of a $d$-regular graph with second eigenvalue bound $\eta$.
For every $Y \in \R^{n \times n}$, if $\hat{Y}$ is the output of the optimization problem 
\eqref{matrix_completion}, then 
\[
	\frac{1}{n^2} \| \hat{Y} - Y \|_F^2 \leq c \gamma_2(Y)^2 \frac{\eta}{d},
\]
where $c = 8 K_G \leq 14.3$ is a universal constant and $\| \cdot \|_F$ is the Frobenius norm.
\end{theorem}

Considering sampling following the biadjacency matrix of a bipartite graph,
we find a similar result which also applies to rectangular matrices. 
If $n=m$ and $d_1 = d_2 = d$, our bound is equivalent to that of Theorem~\ref{thm:square_MSE},
but with constants improved by a factor of two due to stronger mixing in bipartite graphs.
Intuitively, using a biadjacency matrix is a ``more random'' way of sampling 
than using an adjacency matrix, since it is not symmetric.

\begin{theorem}
\label{thm:rect_MSE}
Let $E$ be the set of edges of a $(d_1, d_2)$-regular graph
with second eigenvalue bound $\eta$.
For every $Y \in \R^{n \times m}$, if $\hat{Y}$ is the output of the optimization problem 
\eqref{matrix_completion}, then 
\[
	\frac{1}{n m} \| \hat{Y} - Y \|_F^2 \leq 
      c \gamma_2(Y)^2 \frac{\eta}{\sqrt{d_1 d_2}} ,
\]
where $c = 4 K_G \leq 7.13$.
\end{theorem}
\begin{proof}
We start by considering a rank-1 sign matrix $S=u v^*$, where $u,v \in \{-1,1 \}^{n \times m}$.
Let $S' = \frac{1}{2} (S + J)$, where $J$ is the all-ones matrix, 
so that $S'$ has the entries of -1 in $S$ replaced by zeros.
Then
$S' = 1_A 1_B^* + 1_{A^c} 1_{B^c}^*$ for subsets 
$A \subset V_1 = [n]$ and $B \subset V_2 = [m]$,
where $A = \{ i: u_i = 1 \} $ and $B = \{ j : v_j = 1 \}$.
Consider the expression
\begin{align*}
\left| \frac{1}{n m} \sum_{i,j} s_{ij} - 
\frac{1}{|E|} \sum_{(i,j) \in E} s_{ij} \right|
&= 
\left| \frac{1}{n m} \sum_{i,j} (2 s'_{ij} - 1) - 
\frac{1}{|E|} \sum_{(i,j) \in E} (2 s'_{ij} - 1) \right| \\
&= 
2 \left| \frac{1}{n m} \sum_{i,j} s'_{ij} - 
\frac{1}{|E|} \sum_{(i,j) \in E} s'_{ij} \right| \\
&= 2 \left| \frac{|A||B| + |A^c| |B^c|}{n m}  - 
\frac{E(A,B) + E(A^c, B^c)}{|E|} \right| \\
&\leq 2 \left| \frac{|A||B|}{nm} - \frac{E(A,B)}{|E|} \right|
+ 
2 \left| \frac{|A^c| |B^c|}{n m}  - 
\frac{E(A^c, B^c)}{|E|} \right|.
\end{align*}
% Using a version of the expander mixing lemma, \citep*{janwa},
% \[
% \left| \frac{E(A,B)}{|E|} - \frac{|A||B|}{n m} \right|
% \leq \frac{\eta}{2 |E|} 
%    \left( |A| + |B| - \frac{|A|^2}{n} - \frac{|B|^2}{m} \right),
% \]
% \todo{does the last inequality really hold??}
% we obtain that
% \begin{align*}
% \left| \frac{1}{n m} \sum_{i,j} s_{ij} - 
% \frac{1}{|E|} \sum_{(i,j) \in E} s_{ij} \right|
% & \leq \frac{\eta}{|E|} 
% \left(  
% |A| + |B| - \frac{|A|^2}{n} - \frac{|B|^2}{m} 
%  +
%  |A^c| + |B^c| - \frac{|A^c|^2}{n} - \frac{|B^c|^2}{m}
% \right)\\
% & \leq \frac{\eta}{|E|} 
% \left( 
% n + m - \frac{|A|^2 + |A^c|^2}{n} - \frac{|B|^2 + |B^c|^2}{m}
% \right) \\
% &\leq \frac{\eta}{|E|} 
% \left( 
% n + m - n (x^2 + (1-x)^2) - m (y^2 + (1-y)^2)
% \right) \\
% &\leq 
% \frac{\eta}{2 |E|} (n + m) = 
% \frac{\eta}{2} \left( \frac{d_1 + d_2}{d_1 d_2} \right), 
% \end{align*}
% since $\frac{1}{2} \leq x^2 + (1-x)^2 \leq 1$ for $x \in [0,1]$.
The following is a bipartite version of the expander mixing lemma 
\citep*{dewinter2012}:
\[
\left| \frac{E(A,B)}{|E|} - \frac{|A||B|}{n m} \right|
\leq 
\frac{\eta}{\sqrt{d_1 d_2}} 
\sqrt{\frac{|A||B|}{nm} 
\left(1- \frac{|A|}{n} \right) \left( 1- \frac{|B|}{m} \right) } \\
= 
\frac{\eta}{\sqrt{d_1 d_2}} 
\sqrt{\frac{|A||B||A^c||B^c|}{(nm)^2}} .
\]
We find that
\begin{align*}
\left| \frac{1}{n m} \sum_{i,j} s_{ij} - 
\frac{1}{|E|} \sum_{(i,j) \in E} s_{ij} \right|
& \leq 
\frac{4 \eta}{\sqrt{d_1 d_2}}
\sqrt{\frac{|A||B||A^c||B^c|}{(nm)^2}}
\\
& = \frac{4 \eta}{\sqrt{d_1 d_2}}
\sqrt{ x y (1-x) (1-y) }\\
& \leq \frac{\eta}{\sqrt{d_1 d_2}},
\end{align*}
since $x y (1-x) (1-y)$ attains a maximal value of $2^{-4}$ 
for $0 \leq x,y \leq 1$.

The rest of the proof develops identical to the results of \citet*{heiman2014},
which we include for completeness.
We apply the result for rank-1 sign matrices to any matrix $R$.
Let $R = \sum_i \alpha_i S^i$, where $S^i$ is a rank-1 sign matrix and $\alpha_i \in \R$.
For a general matrix $R$, this might require many rank-1 sign matrices.
Define the sign nuclear norm $\nu (R) = \sum_i |\alpha_i|$. 
Then,
\[
\left| \frac{1}{n m} \sum_{i,j} r_{ij} - 
\frac{1}{|E|} \sum_{(i,j) \in E} r_{ij} \right|
\leq
\nu(R) \frac{\eta}{\sqrt{d_1 d_2}} .
\]
It is a consequence of Grothendieck's inequality,
a well-known theorem in functional analysis,
that there exists a universal constant 
$1.5 \leq K_G \leq 1.8$ 
so that
$\gamma_2 (X) \leq \nu(X) \leq K_G \gamma_2(X)$
for any real matrix $X$; see \citet*{heiman2014}.

Now, let the matrix of residuals 
$R = (\hat{Y}-Y) \circ (\hat{Y} - Y)$, where $\circ$ is the 
Hadamard entry-wise product of two matrices,
so that $R_{ij} = (\hat{Y}_{ij} - Y_{ij})^2$.
Since 
\[
\frac{1}{|E|} \sum_{(i,j) \in E} r_{ij} = 0 ,
\]
we conclude that 
\[
\frac{1}{n m} \sum_{i,j} r_{ij} 
\leq
\nu(R) \frac{\eta}{\sqrt{d_1 d_2}} 
\leq 
K_G \gamma_2(R) \frac{\eta}{\sqrt{d_1 d_2}}.
\]
Furthermore, $\gamma_2 (R) \leq \gamma_2 (\hat{Y}-Y)^2 \leq ( \gamma_2(\hat{Y}) + \gamma_2(Y) ) ^2 $
by \eqref{eq:submultiplicative} and the triangle inequality.
Since $\hat{Y}$ is the output of the algorithm and $Y$ is a feasible solution,
$\gamma_2 (\hat{Y}) \leq \gamma_2 (Y)$.
Thus, $\gamma_2(R) \leq 4 \gamma_2(Y)^2$ and the proof is finished.
\end{proof}

% \begin{remark}
% If we minimize the trace norm of the solution, 
% which is a more practical method than working with $\gamma_2$,
% the same bounds hold in terms of $\| Y \|_{\rm Tr}$.
% This is because $\gamma_2 (Y) \leq \| Y \|_{\rm Tr}$. 
% We only need to modify the final part of the proof.
% \end{remark}

\subsection{Noisy matrix completion bounds}

Furthermore, our analysis easily extends to the case where
the matrix we observe is corrupted with noise. 
As mentioned in the above remark, similar results will hold for the trace norm.
In the noisy case, we solve the problem
\begin{equation}
\label{matrix_completion_2}
\begin{aligned}
& \underset{X}{\text{minimize}}
& & \gamma_2(X) \\
& \text{subject to}
& & \frac{1}{|E|} \sum_{(i,j) \in E} ( X_{ij} - Z_{ij})^2 \leq \delta^2
\end{aligned}
\end{equation}
and obtain the following theorem:
\begin{theorem}
Suppose we observe $Z_{ij} = Y_{ij} + \epsilon_{ij}$ with bounded error
\[
 \frac{1}{|E|} \sum_{(i,j) \in E} \epsilon_{ij}^2 \leq \delta^2 .
\]
Then solving the optimization problem \eqref{matrix_completion_2}
will yield a bound of
\[
\frac{1}{n m} \| \hat{Y} - Y \|_F^2 \leq 
      c \gamma_2(Y)^2 \frac{\eta}{\sqrt{d_1 d_2}} + 4 \delta^2,
\]
where $c = 4 K_G \leq 7.13$.
\end{theorem}
\begin{proof}
Denote $\hat{Y}$ the solution to 
P\ref{matrix_completion_2}.
It will be useful to introduce the sampling operator
$\mathcal{P}_E : \R^{n \times m} \to \R^{n \times m}$,
where
$(\mathcal{P}_E (X) )_{ij} = X_{ij}$
if $(i,j) \in E$ and 0 otherwise.
Again let $R = (\hat{Y} - Y) \circ (\hat{Y} - Y)$
be the matrix of squared errors,
then 
\begin{align*}
\left| 
\frac{1}{nm} \| \hat{Y} - Y \|_F^2 - 
\frac{1}{|E|} \| \mathcal{P}_E (\hat{Y} - Y) \|_F^2 
\right| 
&=
\left| \frac{1}{n m} \sum_{i,j} (\hat{Y}_{ij} - Y_{ij})^2 - 
\frac{1}{|E|} \sum_{(i,j) \in E} (\hat{Y}_{ij} - Y_{ij})^2 \right| \\
& \leq K_G \gamma_2(R) \frac{\eta}{\sqrt{d_1 d_2}}.
\end{align*}
However, since $Y$ is a feasible solution to
P\ref{matrix_completion_2}, we have
\[
\gamma_2 (\hat{Y}) \leq \gamma_2 (Y).
\]
Applying \eqref{eq:submultiplicative} and the triangle inequality, 
\[
\gamma_2 (R) \leq \left( \gamma_2(\hat{Y} - Y) \right)^2
\leq \left( \gamma_2(\hat{Y}) + \gamma_2(Y) \right)^2
\leq 4 \gamma_2(Y)^2 .
\]
Using the triangle inequality again gives
\[
\| \mathcal{P}_E (\hat{Y} - Y) \|_F
\leq
\| \mathcal{P}_E (\hat{Y} - Z) \|_F
+
\| \mathcal{P}_E (Z - Y) \|_F
\leq 2 \delta \sqrt{|E|},
\]
taking into account the bound on the observation errors.
Because
\[
\frac{1}{nm} \| \hat{Y} - Y \|_F^2
\leq
\left| 
 \frac{1}{nm} \| \hat{Y} - Y \|_F^2 - 
 \frac{1}{|E|} \| \mathcal{P}_E (\hat{Y} - Y) \|_F^2 
\right| 
+
\frac{1}{|E|} \| \mathcal{P}_E (\hat{Y} - Y) \|_F^2 
\]
we get the final bound
\[
\frac{1}{n m} \| \hat{Y} - Y \|_F^2
\leq
4 K_G \gamma_2(Y)^2 \frac{\eta}{\sqrt{d_1 d_2}} + 4 \delta^2.
\]
\end{proof}

\subsection{Application of the spectral gap}

Theorem~\ref{thm:rect_MSE} provides a bound on the mean squared error
of the approximation $X$. 
Directly applying Theorem~\ref{cor:gap_A}, 
we obtain the following bound on 
the generalization error of the algorithm using a 
random biregular, bipartite graph:
\begin{corollary}
Let $E$ be sampled from a $\mathcal{G}(n,m, d_1, d_2)$ random graph.
For every $Y \in \R^{n \times m}$, if $\hat{Y}$ is the output of the optimization problem 
\eqref{matrix_completion}, then 
\[
	\frac{1}{n m} \| \hat{Y} - Y \|_F^2 \leq 
      c \gamma_2(Y)^2 \frac{\sqrt{d_1 - 1} + \sqrt{d_2 - 1} + \epsilon_n }{\sqrt{d_1 d_2}} ,
\]
where $c = 4 K_G \leq 7.13$ is a universal constant.
\end{corollary}

